\section{适用范围与局限}

本文的机制表述适用于测试图像内部存在可估计正常区域的局部异常定位场景。在异常覆盖范围很大、整幅图像发生分布偏移、或正常区域本身极难识别时，图像条件化参照会受到异常证据影响。当前方法通过保守正常侧校准控制这一影响；整图异常和正常区域难以分离的场景更适合结合全局判别机制。

其次，本文强调局部响应的语境依赖性；不同类别中的语境来自材质纹理、结构边界、光照变化或成像噪声等因素。现有 boundary-aware reliability 主要处理“高频纹理与结构边界混淆”这一类问题；逻辑异常、几何关系异常或全局语义异常对应更强的全局关系建模。

再次，本文新增模块在测试时保持模型参数固定；底座 CLIP 与 AnomalyCLIP 表征及提示向量来自大规模预训练或辅助异常数据。因而，本文结论聚焦于“在给定冻结表征之上，新增校准模块如何使用单图证据”。

最后，本文的机制性结论来自 MVTec/VisA 受控消融。五数据集系统级结果展示整体表现，其中 MPDD、BTAD 和 DTD-Synthetic 包含 global setting 与 dataset-tuned upper bound 的差异，适合在实验节或附录中单独标注。进一步扩展外部 baseline 或跨数据集比较时，split、backbone、input resolution、post-processing、训练数据和 metric implementation 将作为实验协议的一部分统一说明；定性图将围绕“局部线索需要图像内参照”的机制案例组织。
